\chapter{Asymmetrische Kryptosysteme}\label{chap:asymmetrische-kryptosysteme}

In den bisher angeschauten Verschlüsselungsverfahren haben wir festgestellt, dass derselbe Schlüssel zur Ver- und Entschlüsselung verwendet wird. Daher spricht man bei diesen Verfahren von \emph{symmetrischen} Verschlüsselungsmethoden. Zudem haben wir gesehen, dass diese entweder unsicher (Caesar, Vigenère) sind, wenn der Schlüssel einfach geknackt werden kann, oder unpraktikabel (\gls{OTP}), wenn der Schlüssel zuerst übertragen werden muss. Diffie \& Hellman kamen daher 1975 auf die Idee, \emph{asymmetrische} Verschlüsselungsverfahren zu erschaffen, welche nach einem anderen Prinzip funktionieren. Im Gegensatz zu symmetrischen Verfahren kann man bei asymmetrischen Verfahren \emph{nicht} von der verschlüsselten Nachricht auf den Schlüssel schliessen, da unterschiedliche Schlüssel zum Ver- und Entschlüsseln verwendet werden.

Die Grundidee hinter asymmetrischen Verschlüsselungsmethoden ist folgende:

Alice (\emoji{woman-technologist}) und Bob (\emoji{man-technologist}) möchten auf verschlüsselte Weise Nachrichten austauschen, die den Anforderungen an sichere Kryptosysteme genügen (siehe \cref{sec:intro}). Bei der asymmetrischen Verschlüsselung generieren sowohl Alice wie Bob jeweils ein Schlüsselpaar, einen sogenannten \textbf{öffentlichen Schlüssel} (\emoji{key}), der von allen Personen gesehen und verwendet werden kann, sowie einen \textbf{privaten Schlüssel} (\emoji{lock}), der nur im Besitz von Alice bzw. Bob ist. Wenn Alice eine Nachricht an Bob senden will, verwendet sie Bobs \emph{öffentlichen} Schlüssel, um die Nachricht zu verschlüsseln. Die Nachricht kann jedoch mit dem öffentlichen Schlüssel nicht entschlüsselt werden, es handelt sich hier sozusagen um eine \emph{Einweg}-Funktion. Stattdessen muss Bob seinen \emph{privaten} Schlüssel zur Entschlüsselung verwenden, also den Schlüssel, auf den \emph{nur} Bob Zugriff hat (siehe \cref{fig:public-key-encryption}). Dies funktioniert in den meisten Fällen auch in die andere Richtung: eine Nachricht, die mit Bobs privatem Schlüssel verschlüsselt worden ist, kann nur mit Bobs öffentlichem Schlüssel entschlüsselt werden. Die Schlüssel sind mathematisch so konstruiert, dass es beinahe unmöglich ist, vom öffentlichen Schlüssel auf den privaten Schlüssel zu schliessen.


Durch den Aufbau asymmetrischer Verschlüsselungsverfahren entfällt die Problematik der Übermittlung des Schlüssels: jede Person generiert ihr eigenes Schlüsselpaar und stellt einen öffentlichen Schlüssel zur Verfügung. Ein Nachteil hierbei ist, dass die Verschlüsselung häufig mathematisch und bezüglich Rechenleistung anspruchsvoller ist, was insbesondere problematisch sein kann bei längeren Nachrichten oder wenn die Antwortzeit minimal sein soll.

\begin{figure}[H]
	\centering
	\begin{tikzpicture}[node distance = 1cm and 2cm, data/.style={rectangle, rounded corners, text width=3cm, minimum width=3cm, minimum height=1cm, text centered, draw=ForestGreen, fill=ForestGreen!30},
			startstop/.style={rectangle, rounded corners, minimum height=1cm, minimum width=3cm, text width=2.5cm, text centered, draw=RoyalBlue, fill=RoyalBlue!30}]

		\footnotesize
		% Nodes
		\node (plaintext) [data] {\emoji{woman-technologist} Alice};
		\node (ciphertext) [data, right=of plaintext] {Verschlüsselter\\Text};
		\node (publickey) [startstop, above=of ciphertext] {\emoji{key} Öffentlicher Schlüssel\\von Bob};

		\node (decryptedtext) [data, below=of ciphertext] {Verschlüsselter\\Text};
		\node (privatekey) [startstop, below=of decryptedtext] {\emoji{lock} Privater Schlüssel\\ von Bob};
		\node (bob) [data, left=of decryptedtext] {\emoji{man-technologist} Bob};


		% Arrows
		\draw [->] (plaintext) -- (ciphertext);
		\draw [->] (ciphertext) -- (decryptedtext);
		\draw [->] (publickey) -- (ciphertext);
		\draw [->] (privatekey) -- (decryptedtext);
		\draw [->] (decryptedtext) -- (bob);

		\node at ($(plaintext.east)!0.5!(ciphertext.west)$) [above] {Klartext};
		\node at ($(publickey.south)!0.5!(ciphertext.north)$) [left] {Verschlüsselung};
		\node at ($(ciphertext.south)!0.5!(decryptedtext.north)$) [left] {Übertragung};
		\node at ($(privatekey.north)!0.5!(decryptedtext.south)$) [left] {Entschlüsselung};
		\node at ($(decryptedtext.west)!0.5!(bob.east)$) [above] {Klartext};
	\end{tikzpicture}
	\caption{Prinzip der Public-Key-Verschlüsselung}
	\label{fig:public-key-encryption}
\end{figure}

\section{\texorpdfstring{\acrshort{RSA}}{RSA}-Verfahren}
\label{subsec:rsa}

Das \acrfull{RSA}-Verfahren ist ein asymmetrisches Kryptosystem, das sowohl für die Verschlüsselung als auch für digitale Signaturen verwendet werden kann. Es basiert auf der Schwierigkeit, grosse Zahlen in ihre Primfaktoren zu zerlegen. Im Folgenden werden die Grundlagen des Verfahrens und ein einfaches Beispiel vorgestellt.

\begin{figure}[H]
	\centering
	\includegraphics[width=0.7\linewidth]{Grundlagen_Info/03_Kryptologie/Figures/rsa.jpg}
	\caption{Ron Rivest, Adi Shamir und Leonard Adleman, die Erfinder des \gls{RSA}-Verfahrens}
	\label{fig:rsa}
\end{figure}


\begin{myprocedure}[\gls{RSA}-Verschlüsselungsverfahren]\label{myprocedure:RSA}
	Das \gls{RSA}-Verfahren besteht aus folgenden Schritten:
\begin{enumerate}
	\item \textbf{Schlüsselerzeugung:}
	      \begin{itemize}
		      \item Wählen Sie \textbf{zwei (möglichst grosse) Primzahlen} \( p \) und \( q \).
		      \item Berechnen Sie das \textbf{\gls{RSA}-Modul} \( n = p \cdot q \).
		      \item Berechnen Sie die \textbf{Eulersche Funktion} \( \varphi(n) = (p-1)(q-1) \). Die Eulersche Funktion \( \varphi(n) \) gibt die Anzahl der positiven ganzen Zahlen kleiner als \( n \) an, die teilerfremd zu \( n \) sind. Für eine Primzahl $p$ gilt \( \varphi(p) = p-1 \), daher \( \varphi(n) = \varphi(p \cdot q) = (p-1)(q-1) \).
		      \item Wählen Sie den \textbf{\emph{Verschlüsselungsexponenten} \( e \)}, so dass dieser teilerfremd zu \( \varphi(n) \) und kleiner als $\varphi(n)$ ist. Dies bedeutet, dass der \gls{GGT} von $e$ und $\varphi(n)$ 1 ist ($\operatorname{ggT}(e, \varphi(n)) = 1$).
		      \item Berechnen Sie den \textbf{\emph{Entschlüsselungsexponenten} \( d \)}, so dass \( (e \cdot d) \equiv 1 \pmod{\varphi(n)} \) (privater Schlüssel). Dies bedeutet: Bei der Division des Produkts $e \cdot d$ durch $\varphi(n)$ bleibt der Rest $1$, sodass $e \cdot d = 1 + k \cdot \varphi(n)$ für eine ganze Zahl $k$ gilt.
		      \item Die Zahlen $p$, $q$ und $\varphi(n)$ werden nun nicht mehr benötigt und können gelöscht werden.
	      \end{itemize}
	\item \textbf{Verschlüsselung:}
	      \begin{itemize}
		      \item Der Absender verschlüsselt eine \textbf{Klartext-Nachricht} \( \mathbf{m} \) (als Zahl) mit dem \textbf{öffentlichen Schlüssel} \( \mathbf{(e, n)} \):
		            \[
			            c = \mathbf{m}^e \mod n
		            \]
		      \item Das Ergebnis \( c \) ist der Kryptotext.
	      \end{itemize}
	\item \textbf{Entschlüsselung:}
	      \begin{itemize}
		      \item Der Empfänger entschlüsselt die \textbf{verschlüsselte Nachricht} \( \mathbf{c} \) mit dem \textbf{privaten Schlüssel} \( \mathbf{(d, n)} \):
		            \[
			            \mathbf{m} = c^d \mod n
		            \]
		      \item Das Ergebnis \( \mathbf{m} \) ist die ursprüngliche Nachricht.
	      \end{itemize}
\end{enumerate}
\end{myprocedure}

\begin{myremark}[Effiziente modulare Exponentiation bei \gls{RSA}]\label{rem:rsa-mod-exp}
	Die Berechnungen $c = m^e \mod n$ und $m = c^d \mod n$ erfordern für praxisnahe Schlüssellängen (z.\,B.\ 2048 Bit) das Potenzieren mit riesigen Zahlen. Wie bereits beim \gls{DHM}-Verfahren (siehe \cref{rem:modular-exp-naive}) kommt auch bei \gls{RSA} die \textbf{modulare Exponentiation} (wiederholtes Quadrieren) zum Einsatz. Dadurch können die Berechnungen in Millisekunden durchgeführt werden, ohne dass astronomisch grosse Zwischenzahlen gespeichert werden müssen.
\end{myremark}

\begin{myexample}
Um \gls{RSA} besser zu verstehen, rechnen wir ein einfaches Beispiel mit kleinen Zahlen durch:
\begin{enumerate}
	\item Wir wählen zwei (für diese Übung kleine) Primzahlen aus, \( p = 3 \) und \( q = 11 \). Daraus folgt:
	      \begin{align*}
		      n & = p \cdot q  \\
		        & = 3 \cdot 11 \\
		        & = 33
	      \end{align*}
	      \begin{align*}
		      \varphi(n) & = (p-1)(q-1) \\
		                 & = 2 \cdot 10 \\
		                 & = 20
	      \end{align*}
	\item Wir wählen \( e = 3 \), da \( \operatorname{ggT}(3, 20) = 1 \).
	\item Wir berechnen \( d \), so dass \( (e \cdot d) \mod \varphi(n) = 1  \) ergibt:
	      \[
		      (d \cdot 3) \mod 20 = 1  \quad \rightarrow \quad d = 7.
	      \]
	\item Der öffentliche Schlüssel ist \( (e, n) = (3, 33) \), der private Schlüssel \( (d, n) = (7, 33) \).
\end{enumerate}

\textbf{Verschlüsselung:} Die Nachricht sei der Klartext \texttt{Code}, welches wir darstellen durch die Position der Buchstaben im Alphabet \( m = 3, 15, 4, 5 \). Berechne für jeden Buchstaben (hier nur für \texttt{C}, also 3, gezeigt):
\begin{align*}
	c & = m^e \mod n  \\
	  & = 3^3 \mod 33 \\
	  & = 27 \mod 33  \\
	  & = 27.
\end{align*}


Der erste Buchstabe des Kryptotext wird also verschlüsselt als \( c = 27\).

\textbf{Entschlüsselung:}
\begin{align*}
	m & = c^d \mod n     \\
	  & = 27^{7} \mod 33 \\
	  & = 3
\end{align*}

Daraus ergibt sich \( m = 3 \), also die ursprüngliche Nachricht.
\end{myexample}

\begin{myexcursion}[Mathematischer Hintergrund: Der Satz von Euler und der \texorpdfstring{\acrshort{RSA}}{RSA}-Korrektheitsbeweis]\label{exc:rsa-beweis}\label{exc:euler}
	Warum stellt die Entschlüsselung $(m^e \bmod n)^d \bmod n$ garantiert die ursprüngliche Nachricht $m$ wieder her? Das mathematische Fundament hierfür liefert ein berühmter Satz des Schweizer Mathematikers Leonhard Euler (1707--1783).

	\textbf{1. Die Eulersche $\varphi$-Funktion:}
	Die Funktion $\varphi(n)$ zählt, wie viele positive ganze Zahlen kleiner als $n$ teilerfremd zu $n$ sind ($\operatorname{ggT}(a, n) = 1$):
	\begin{itemize}
		\item Für eine Primzahl $p$ sind alle Zahlen von $1$ bis $p-1$ teilerfremd zu $p$:
		      \[
			      \varphi(p) = p - 1.
		      \]
		\item Für das \gls{RSA}-Modul $n = p \cdot q$ (Produkt zweier verschiedener Primzahlen) gilt:
		      \[
			      \varphi(n) = \varphi(p \cdot q) = (p - 1)(q - 1).
		      \]
	\end{itemize}

	\textbf{2. Der Satz von Euler:}
	Ist eine Basis $a$ teilerfremd zu $n$ ($\operatorname{ggT}(a, n) = 1$), so besagt der \textbf{Satz von Euler}:
	\[
		a^{\varphi(n)} \equiv 1 \pmod n.
	\]
	\textbf{Beispiel:} Wählen wir $n = 10 = 2 \cdot 5$, so ist $\varphi(10) = (2-1)(5-1) = 4$ (die $4$ teilerfremden Zahlen sind $\{1, 3, 7, 9\}$). Für die Basis $a = 3$ gilt tatsächlich:
	\[
		3^{\varphi(10)} = 3^4 = 81 \equiv 1 \pmod{10}.
	\]

	\textbf{3. Schrittweiser Beweis der \gls{RSA}-Entschlüsselung:}
	Bei der Schlüsselerzeugung wird $d$ so gewählt, dass $e \cdot d \equiv 1 \pmod{\varphi(n)}$ gilt. Das bedeutet: Bei der Division von $e \cdot d$ durch $\varphi(n)$ bleibt der Rest $1$, also existiert eine ganze Zahl $k$ mit
	\[
		e \cdot d = 1 + k \cdot \varphi(n) = k \cdot \varphi(n) + 1.
	\]
	Daraus folgt für die Entschlüsselung eines Kryptotexts $c = m^e \bmod n$:
	\begin{align*}
		\text{Entschlüsselter Text} 
		  & = c^d \bmod n                                          && (\text{Definition der Entschlüsselung}) \\[0.4em]
		  & = (m^e \bmod n)^d \bmod n                              && (\text{Kryptotext } c = m^e \bmod n \text{ einsetzen}) \\[0.4em]
		  & \equiv (m^e)^d \pmod n                                 && (\text{Potenzierungsregel, s.~\cref{exc:potenzierungsregel}}) \\[0.4em]
		  & \equiv m^{e \cdot d} \pmod n                           && (\text{Potenzgesetz: } (a^b)^c = a^{b \cdot c}) \\[0.4em]
		  & \equiv m^{k \cdot \varphi(n) + 1} \pmod n              && (\text{Schlüsselbedingung: } e \cdot d = 1 + k \cdot \varphi(n)) \\[0.4em]
		  & \equiv \left(m^{\varphi(n)}\right)^k \cdot m^1 \pmod n && (\text{Potenzgesetz: } a^{x+y} = a^x \cdot a^y) \\[0.4em]
		  & \equiv (1)^k \cdot m \pmod n                           && (\text{Satz von Euler: } m^{\varphi(n)} \equiv 1) \\[0.4em]
		  & \equiv 1 \cdot m \pmod n                               && (1^k = 1) \\[0.4em]
		  & = \mathbf{m}                                           && (\text{Klartext wiederhergestellt!})
	\end{align*}
	
	\textbf{Zusammenfassung der Kernidee:} Durch die Wahl von $e \cdot d = 1 + k \cdot \varphi(n)$ durchläuft der Exponent genau $k$ volle Umläufe der Länge $\varphi(n)$ modulo $n$, die nach dem Satz von Euler jeweils neutral zu $1$ werden ($1^k = 1$). Es bleibt exakt der Faktor $m^1 = m$ übrig.
\end{myexcursion}

\begin{myexercise}\label{myexercise:rsa}
	Alice möchte eine Nachricht an Bob mit \gls{RSA} verschlüsseln. Bob wählt die Hilfs-Primzahlen \( p = 5 \) und \( q = 7 \). Generieren Sie den öffentlichen Schlüssel ($e, n$) und den privaten Schlüssel ($d, n$) von Bob, indem Sie folgende Schritte ausführen:
	\begin{enumerate}
		\item Berechnen Sie \( n \) und \( \varphi(n) \).
		\item Sie wählen $e$ und $d$ aus.
	\end{enumerate}

	Verschlüsseln Sie nun die Nachricht \( m = 9 \) mit dem öffentlichen Schlüssel \( (e, n) \).

	Entschlüsseln Sie danach die verschlüsselte Nachricht \( c \) wieder mit dem privaten Schlüssel \( (d, n) \).
	\begin{myanswer}
		\begin{enumerate}
			\item \( n = p \cdot q = 5 \cdot 7 = 35 \), \( \varphi(n) = (p-1)(q-1) = 4 \cdot 6 = 24 \).
			\item Für $e$ können wir beispielsweise 5 wählen, da 5 teilerfremd mit 24 ist.
			\item \( (5 \cdot d) \mod 24 = 1\). Dies finden wir beispielsweise mit \( d = 5 \).
			\item Verschlüsselung:
			      \begin{align*}
				      c & = m^e\mod n   \\
				        & = 9^5 \mod 35 \\
				        & = 4
			      \end{align*}.
			\item Entschlüsselung:
			      \begin{align*}
				      m & = c^d\mod n    \\
				        & =4^{5} \mod 35 \\
				        & = 9
			      \end{align*}
		\end{enumerate}
	\end{myanswer}
\end{myexercise}

\begin{myexercise}
	Ist die Wahl von \( p \) und \( q \) im obigen Beispiel sinnvoll? Begründen Sie Ihre Antwort.
	\begin{myanswer}
		Nein, die Wahl von \( p \) und \( q \) ist nicht sinnvoll, da beide Primzahlen zu klein sind. In der Praxis sollten \( p \) und \( q \) jeweils mindestens 2048 Bit lang sein, um eine ausreichende Sicherheit zu gewährleisten. Ansonsten kann es, wie im vorliegenden Beispiel, vorkommen, dass der private Schlüssel leicht durch Faktorisierung von \( n \) berechnet werden kann, oder dass der private Schlüssel sogar gleich dem öffentlichen Schlüssel ist (wie in diesem Beispiel, wo \( d = e = 5 \)).
	\end{myanswer}
\end{myexercise}

\begin{mychallenge}\label{mychallenge:rsa-pair}
 Erstellen Sie nun zu zweit jeweils ein Schlüsselpaar mit dem \gls{RSA}-Verfahren. Verwenden Sie dazu Primzahlen \( p \) und \( q \) mit jeweils mindestens 2 Ziffern. Tauschen Sie danach Ihre öffentlichen Schlüssel aus und verschlüsseln Sie eine Nachricht (eine Zahl) an Ihren Partner. Entschlüsseln Sie danach die Nachricht wieder.
 
 Eine Liste der ersten 1000 Primzahlen finden Sie hier: \url{https://en.wikipedia.org/wiki/List_of_prime_numbers}.
\end{mychallenge}

\begin{mychallenge}
Verwenden Sie Ihre \gls{RSA}-Schlüssel aus \ref{mychallenge:rsa-pair}, um eine echte Nachricht als Zahl auszutauschen (ein Wort). Um ein Wort in Zahlen umzuwandeln (Buchstabe für Buchstabe) können Sie folgendes Python-Skript verwenden:
\begin{lstlisting}
def text_to_numbers(text):
	numbers = []
	for char in text.upper():
		if char.isalpha():  # Nur Buchstaben berücksichtigen
			numbers.append(ord(char) - ord('A') + 1)  # A=1, B=2, ..., Z=26
	return numbers

print(text_to_numbers("ABC"))  # Ausgabe: [1, 2, 3]
\end{lstlisting}
\end{mychallenge}

\begin{mychallenge}
	Wählen Sie zwei Primzahlen \( p \) und \( q \) mit jeweils mindestens 3 Ziffern. Generieren Sie den öffentlichen und privaten Schlüssel. Verschlüsseln Sie eine Nachricht Ihrer Wahl und entschlüsseln Sie diese wieder. Verwenden Sie Python, um die Berechnungen durchzuführen.

	Die Zahl $d$ kann in Python folgendermassen berechnet werden:

	Eine Liste der ersten 1000 Primzahlen finden Sie hier: \url{https://en.wikipedia.org/wiki/List_of_prime_numbers}
	\begin{lstlisting}
d = pow(e, -1, phi_n)
	\end{lstlisting}
	\begin{myanswer}
		Beispiel mit \( p = 101 \) und \( q = 113 \):
		\lstinputlisting{Grundlagen_Info/03_Kryptologie/Python package/encryption/rsa_simple.py}
	\end{myanswer}
\end{mychallenge}

\subsection*{Sicherheit des \texorpdfstring{\acrshort{RSA}}{RSA}-Verfahrens}
Die Sicherheit der reinen \gls{RSA}-Verschlüsselung beruht auf der mathematischen Schwierigkeit der \textbf{Primfaktorzerlegung} (Faktorisierungsproblem). Ein Angreifer (Eve) hat Zugriff auf den öffentlichen Schlüssel $(e, n)$ sowie den abgefangenen Kryptotext $c$. Um daraus den privaten Entschlüsselungsexponenten $d$ zu berechnen, müsste Eve den Wert der Eulerschen Phi-Funktion $\varphi(n) = (p-1)(q-1)$ kennen. Dazu müsste sie die Zahl $n$ in ihre ursprünglichen Primfaktoren $p$ und $q$ zerlegen ($n = p \cdot q$).

Während die Multiplikation zweier grosser Primzahlen für moderne Computer im Bruchteil einer Millisekunde durchführbar ist, existiert bis heute kein bekanntes effizientes Verfahren, um ein Produkt aus zwei sehr grossen Primzahlen wieder zu faktorisieren. Werden ausreichend lange Primzahlen gewählt (in der Praxis mindestens 2048 bis 4096 Bit, d.\,h.\ Zahlen mit über 600 Dezimalstellen), würde ein Kryptoanalyst selbst mit den leistungsfähigsten Supercomputern Millionen von Jahren für die Faktorisierung benötigen.

\paragraph*{Sicherheitsproblem: Fehlende Authentizität des Absenders}
Obwohl das \gls{RSA}-Verfahren die \textbf{Vertraulichkeit} der Nachricht schützt (nur der Empfänger Bob kann mit seinem privaten Schlüssel $d$ entschlüsseln), bietet reine Verschlüsselung \textbf{keine Authentizität}: Da Bobs öffentlicher Schlüssel $(e, n)$ öffentlich bekannt ist, kann \emph{jede} beliebige Person (also auch ein Angreifer wie Eve) eine Nachricht verschlüsseln und an Bob schicken, während sie behauptet, Alice zu sein. Bob hat keine Möglichkeit, anhand der verschlüsselten Nachricht nachzuweisen, wer sie tatsächlich verfasst hat.

\paragraph*{Lösung: Digitale Signaturen durch Rollenumkehr der Schlüssel}
Um diese Sicherheitslücke zu schliessen, nutzt man die mathematische Symmetrie des \gls{RSA}-Verfahrens: Die Rollen des privaten und des öffentlichen Schlüssels lassen sich vertauschen. Da die Multiplikation im Exponenten kommutativ ist ($d \cdot e = e \cdot d$), gilt nach dem Beweis aus \cref{exc:rsa-beweis} genauso die umgekehrte Reihenfolge:
\[
	(m^d)^e = m^{d \cdot e} = m^{e \cdot d} \equiv m \pmod n.
\]
Eine Person kann eine Nachricht also mit ihrem \emph{privaten} Schlüssel $d$ signieren ($s = m^d \bmod n$). Anschliessend kann jede andere Person mit dem zugehörigen \emph{öffentlichen} Schlüssel $e$ überprüfen ($s^e \equiv m \pmod n$), ob die Nachricht tatsächlich vom Schlüsselinhaber stammt und unverändert geblieben ist. Dies bildet das Fundament für \textbf{digitale Signaturen}.

\subsection{Digitale Signaturen mit \texorpdfstring{\acrshort{RSA}}{RSA}}
Neben der Verschlüsselung von Nachrichten kann das \gls{RSA}-Verfahren auch zur Erstellung digitaler Signaturen verwendet werden. Digitale Signaturen dienen dazu, die Authentizität und Integrität einer Nachricht zu gewährleisten. Hierbei wird die Nachricht mit dem privaten Schlüssel des Absenders signiert, sodass der Empfänger die Signatur mit dem öffentlichen Schlüssel des Absenders überprüfen kann. In diesem Abschnitt wird das Verfahren zur Erstellung und Überprüfung digitaler Signaturen mit \gls{RSA} erläutert. Dabei wird in einem ersten Schritt angenommen, dass die Nachricht $m$ unverschlüsselt übertragen wird -- das Ziel dabei ist, dass der Empfänger die Nachricht eindeutig der Absenderin zuordnen und gleichzeitig überprüfen kann, dass diese nicht verändert wurde.

\begin{myprocedure}[Digitale Signaturen mit \gls{RSA}]
	Die Erstellung und Überprüfung digitaler Signaturen mit dem \gls{RSA}-Verfahren erfolgt in folgenden Schritten:
	\begin{enumerate}
		\item \textbf{Signaturerstellung (durch die Absenderin):}
		      \begin{itemize}
			      \item Alice besitzt die Nachricht (als Zahl) $m$ mit $m < n$.
			      \item Sie berechnet die Signatur $s$ mit ihrem \textbf{privaten Schlüssel} $(d, n)$:
			            \[
				            s = m^d \bmod n
			            \]
			      \item Das Ergebnis $s$ ist die digitale Signatur. Alice sendet die Nachricht $m$ und die Signatur $s$ als Paar $(m, s)$ an Bob.
		      \end{itemize}
		\item \textbf{Signaturüberprüfung (durch den Empfänger):}
		      \begin{itemize}
			      \item Bob erhält die Nachricht $m$ und die Signatur $s$.
			      \item Er entschlüsselt die Signatur mit Alices \textbf{öffentlichem Schlüssel} $(e, n)$:
			            \[
				            m' = s^e \bmod n
			            \]
			      \item Wenn $m' = m$, ist die Signatur gültig (die Nachricht stammt garantiert von Alice und wurde nicht verändert). Andernfalls ist sie ungültig.
		      \end{itemize}
	\end{enumerate}
\end{myprocedure}

\begin{myexample}
	Angenommen, Alice möchte eine Nachricht $m = 42$ signieren. Sie will die (unverschlüsselte) Nachricht 42 also verschicken -- dabei soll eindeutig überprüfbar sein, dass Alice die Absenderin ist. Sie verwendet ihren privaten Schlüssel $(d, n) = (9677, 12317)$ und den öffentlichen Schlüssel $(e, n) = (5, 12317)$.

	\textbf{Signaturerstellung:}
	\[
		s = m^d \bmod n = 42^{9677} \bmod 12317 = 161
	\]
	Die digitale Signatur ist also $s = 161$.

	\textbf{Signaturüberprüfung:}
	\[
		m' = s^e \bmod n = 161^5 \bmod 12317 = 42
	\]
	Da $m' = m = 42$, ist die Signatur gültig.
\end{myexample}

\begin{myexercise}
	Alice möchte die Nachricht $m = 15$ signieren. Ihr privater Schlüssel ist $(d, n) = (7, 33)$ und ihr öffentlicher Schlüssel ist $(e, n) = (3, 33)$.

	Erstellen Sie die digitale Signatur $s$ für die Nachricht $m$.

	Überprüfen Sie danach die Signatur mit dem öffentlichen Schlüssel.
	\begin{myanswer}
		\textbf{Signaturerstellung:}
		\[
			s = m^d \bmod n = 15^7 \bmod 33 = 27
		\]
		Die digitale Signatur ist also $s = 27$.

		\textbf{Signaturüberprüfung:}
		\[
			m' = s^e \bmod n = 27^3 \bmod 33 = 15
		\]
		Da $m' = m = 15$, ist die Signatur gültig.
	\end{myanswer}
\end{myexercise}

\begin{myremark}[Praxis-Exkurs: Digitale Fingerabdrücke mit Hashfunktionen]
	In der Praxis ist eine Nachricht $m$ (z.\,B.\ ein \gls{PDF}-Dokument, ein Foto oder ein Software-Update) oft viele Megabytes gross. Da \gls{RSA}-Berechnungen mit riesigen Zahlen extrem rechenintensiv sind und die Nachricht $m$ kleiner als der \gls{RSA}-Modul $n$ sein muss, signiert man in realen Systemen nicht die gesamte Datei direkt, sondern ihren \textbf{kryptographischen Hashwert}:
	\begin{itemize}
		\item \textbf{Was ist eine Hashfunktion?} Eine mathematische Einwegfunktion (z.\,B.\ \gls{SHA}-256), die aus einer beliebig langen Eingabedatei $m$ einen kurzen Prüfwert $H(m)$ fester Länge (den sogenannten \textbf{Hashwert} bzw.\ \enquote{digitalen Fingerabdruck}, z.\,B.\ 256 Bit) berechnet. 
		\item \textbf{Zentrale Eigenschaft (Lawineneffekt):} Bereits die Änderung eines einzigen Zeichens im Dokument führt zu einem vollkommen anderen Hashwert. Zudem ist es praktisch unmöglich, zwei verschiedene Dokumente mit demselben Hashwert zu finden (\emph{Kollisionsresistenz}).
		\item \textbf{Ablauf in der Praxis:} Alice berechnet $h = H(m)$ und signiert nur diesen kurzen Fingerabdruck: $s = h^d \bmod n$. Bob berechnet seinerseits $H(m)$ aus der erhaltenen Datei und vergleicht es mit $h' = s^e \bmod n$.
	\end{itemize}
\end{myremark}

\subsection*{Klartext-Signatur vs.\ Signatur mit Verschlüsselung}
In der Kryptographie ist es entscheidend, zwischen den verschiedenen \textbf{Schutzzielen} zu unterscheiden:
\begin{itemize}
	\item \textbf{Verschlüsselung} schützt die \textbf{Vertraulichkeit} (\emph{Wer darf die Nachricht lesen?}).
	\item \textbf{Digitale Signaturen} schützen die \textbf{Authentizität} (\emph{Wer ist der Urheber?}) und die \textbf{Integrität} (\emph{Wurde die Nachricht manipuliert?}).
\end{itemize}

\paragraph*{Fall 1: Reine digitale Signatur (Klartextübertragung)}
Bei einer reinen Signatur wird die Nachricht $m$ \textbf{im Klartext} zusammen mit der Signatur $s$ übertragen.
\begin{itemize}
	\item \textbf{Warum Klartext?} Eine Signatur soll beweisen, von wem das Dokument stammt und dass nichts verändert wurde -- sie dient \emph{nicht} der Geheimhaltung. Zur Überprüfung ($m' = s^e \bmod n$) muss dem Empfänger der Klartext $m$ zwingend vorliegen, um den Vergleich $m' \stackrel{?}{=} m$ durchführen zu können.
	\item \textbf{Typische Einsatzgebiete:}
	      \begin{itemize}
		      \item \textbf{Software- und Betriebssystem-Updates:} Das Update soll von allen heruntergeladen werden können, die Signatur garantiert jedoch, dass es unverändert vom echten Hersteller stammt.
		      \item \textbf{Gesetze, Behördendokumente \& Verträge:} Der Inhalt ist öffentlich bestimmt, die Signatur garantiert die Rechtsgültigkeit (z.\,B.\ digital signierte \glspl{PDF}).
		      \item \textbf{Git-Commits:} Quellcode ist einsehbar, aber Commits werden kryptographisch signiert.
		      \item \textbf{Digitale Zertifikate (\gls{PKI}):} Zertifikate müssen für jedermann lesbar sein.
	      \end{itemize}
\end{itemize}

\paragraph*{Fall 2: Kombinierte Signatur \& Verschlüsselung (\emph{Sign-then-Encrypt})}
Soll eine Nachricht sowohl \textbf{geheim} als auch \textbf{authentisch} sein, kombiniert man beide Verfahren (siehe \cref{fig:sign-then-encrypt-detail}):
\begin{enumerate}
	\item Alice signiert den Klartext $m$ mit ihrem privaten Schlüssel $d_A$: $s = m^{d_A} \bmod n_A$.
	\item Alice verschlüsselt das Paket $(m, s)$ mit Bobs öffentlichem Schlüssel $e_B$: $c = \text{Enc}_{e_B}(m, s)$.
	\item Über den Kanal wird ausschliesslich das Chiffrat $c$ übertragen (für Dritte unlesbar).
	\item Bob entschlüsselt $c$ mit seinem privaten Schlüssel $d_B \rightarrow (m, s)$.
	\item Bob prüft die Signatur $s$ mit Alices öffentlichem Schlüssel $e_A \rightarrow s^{e_A} \bmod n_A \stackrel{?}{=} m$.
\end{enumerate}
\begin{itemize}
	\item \textbf{Typische Einsatzgebiete:}
	      \begin{itemize}
		      \item \textbf{Vertrauliche E-Mails:} Verschlüsselte und signierte Nachrichten via \gls{SMIME} oder \gls{PGP}/\gls{GPG}.
		      \item \textbf{Ende-zu-Ende-verschlüsselte Messenger:} Signal, WhatsApp oder Threema.
		      \item \textbf{Geschäftsgeheimnisse \& Patientendaten:} Vertrauliche Verträge und medizinische Akten.
	      \end{itemize}
\end{itemize}

\begin{figure}[H]
	\centering
	\begin{tikzpicture}[
		scale=0.80, every node/.style={transform shape},
		actor/.style={rectangle, rounded corners, draw=black, fill=gray!15, minimum width=3.4cm, minimum height=1.1cm, font=\bfseries, align=center, inner sep=4pt},
		stepbox/.style={rectangle, rounded corners, draw=#1, fill=#1!8, line width=0.8pt, font=\small, align=center, inner sep=6pt},
		msgarrow/.style={->, >=stealth, line width=1pt, #1},
		badge/.style={circle, fill=#1, text=white, font=\bfseries\footnotesize, inner sep=2pt}
	]
		% Spaltenkoordinaten
		\def\colA{0}   % Alice
		\def\colB{9.5} % Bob

		% --- AKTEURE (Kopfzeile) ---
		\node[actor] (alice) at (\colA, 0) {Alice\\{\footnotesize (Absenderin)}\\[2pt]{\scriptsize priv.: $d_A$, Bobs öff.: $(e_B, n_B)$}};
		\node[actor] (bob) at (\colB, 0) {Bob\\{\footnotesize (Empfänger)}\\[2pt]{\scriptsize priv.: $d_B$, Alices öff.: $(e_A, n_A)$}};

		% Vertikale Hilfslinien
		\draw[dashed, gray!40] (alice.south) -- (\colA, -10.2);
		\draw[dashed, gray!40] (bob.south) -- (\colB, -10.2);

		% --- SCHRITT 1 & 2 BEI ALICE ---
		\node[font=\footnotesize\bfseries\color{RoyalBlue!80!black}, fill=white, inner sep=2pt] at (\colA, -1.1) {Phase 1: Signieren \& Verschlüsseln};

		\node[stepbox=RoyalBlue] (alice_proc) at (\colA, -2.7) {
			\textbf{1. Signatur erzeugen (mit $d_A$):}\\
			$s := m^{d_A} \bmod n_A$\\[0.4em]
			\textbf{2. Verschlüsseln (mit $e_B$):}\\
			$c_m := m^{e_B} \bmod n_B$\\
			$c_s := s^{e_B} \bmod n_B$
		};

		% --- NACHRICHTENÜBERTRAGUNG ---
		\node[badge=DarkOrchid] (b_msg) at (\colA+0.6, -5.1) {\faLock};
		\draw[msgarrow=DarkOrchid] (b_msg.east) -- (\colB-0.6, -5.1)
			node[midway, above=3pt, font=\footnotesize\bfseries] {Chiffrate $(c_m, c_s)$}
			node[midway, below=3pt, font=\scriptsize\color{gray!70!black}] {(vertraulich \& verschlüsselt: Eve kann weder $m$ noch $s$ lesen)};

		% --- SCHRITT 3 & 4 BEI BOB ---
		\node[font=\footnotesize\bfseries\color{ForestGreen!80!black}, fill=white, inner sep=2pt] at (\colB, -6.8) {Phase 2: Entschlüsseln \& Verifizieren};

		\node[stepbox=ForestGreen] (bob_proc) at (\colB, -8.6) {
			\textbf{3. Entschlüsseln (mit $d_B$):}\\
			$m := (c_m)^{d_B} \bmod n_B$\\
			$s := (c_s)^{d_B} \bmod n_B$\\[0.4em]
			\textbf{4. Signatur prüfen (mit $e_A$):}\\
			$m' := s^{e_A} \bmod n_A$\\
			$\rightarrow \text{Gilt } m' = m\text{?} \implies \text{Echt, unverfälscht \& geheim!} \textcolor{ForestGreen}{\mathbf{\checkmark}}$
		};
	\end{tikzpicture}
	\caption{Detaillierter Ablauf von \emph{Sign-then-Encrypt}: Kombinierte digitale Signatur mit Verschlüsselung}
	\label{fig:sign-then-encrypt-detail}
\end{figure}

\begin{myexercise}[Kombinierte Signatur und Verschlüsselung in der Praxis]
	Alice möchte die geheime Botschaft $m = 4$ an Bob senden. Die Übertragung soll sowohl \textbf{vertraulich} (geheim) als auch \textbf{authentisch} (nachweisbar von Alice) sein. Dazu nutzen beide das \gls{RSA}-Verfahren für \emph{Sign-then-Encrypt}.

	Folgende Schlüsselpaare sind bekannt:
	\begin{itemize}
		\item \textbf{Alice:} Öffentlicher Schlüssel $(e_A, n_A) = (3, 33)$, privater Schlüssel $d_A = 7$.
		\item \textbf{Bob:} Öffentlicher Schlüssel $(e_B, n_B) = (3, 55)$, privater Schlüssel $d_B = 27$.
	\end{itemize}

	\begin{enumerate}
		\item \textbf{Signaturerstellung:} Berechnen Sie die digitale Signatur $s$ von Alice für die Nachricht $m = 4$.
		\item \textbf{Verschlüsselung:} Berechnen Sie die beiden Chiffrate $c_m$ (verschlüsselte Nachricht) und $c_s$ (verschlüsselte Signatur), die Alice an Bob sendet.
		\item \textbf{Entschlüsselung durch Bob:} Berechnen Sie, wie Bob aus den empfangenen Chiffraten $(c_m, c_s)$ den Klartext $m$ und die Signatur $s$ wiederherstellt.
		\item \textbf{Signaturprüfung durch Bob:} Wie verifiziert Bob anhand von $s$ und Alices öffentlichem Schlüssel, dass die Nachricht tatsächlich von Alice stammt?
	\end{enumerate}
	\begin{myanswer}
		\begin{enumerate}
			\item \textbf{Signaturerstellung (mit $d_A = 7$ und $n_A = 33$):}
			      \[
				      s = m^{d_A} \bmod n_A = 4^7 \bmod 33.
			      \]
			      Wegen $4^3 = 64 \equiv 31 \equiv -2 \pmod{33}$ gilt:
			      \[
				      4^7 = (4^3)^2 \cdot 4 \equiv (-2)^2 \cdot 4 = 4 \cdot 4 = 16 \pmod{33} \implies s = 16.
			      \]
			\item \textbf{Verschlüsselung für Bob (mit $e_B = 3$ und $n_B = 55$):}
			      \begin{align*}
				      c_m & = m^{e_B} \bmod n_B = 4^3 \bmod 55 = 64 \bmod 55 = 9, \\
				      c_s & = s^{e_B} \bmod n_B = 16^3 \bmod 55 = 4096 \bmod 55 = 26.
			      \end{align*}
			      Alice überträgt das Chiffratpaar $(c_m, c_s) = (9, 26)$.
			\item \textbf{Entschlüsselung durch Bob (mit $d_B = 27$ und $n_B = 55$):}
			      \begin{align*}
				      m & = (c_m)^{d_B} \bmod n_B = 9^{27} \bmod 55 = 4, \\
				      s & = (c_s)^{d_B} \bmod n_B = 26^{27} \bmod 55 = 16.
			      \end{align*}
			      (Zwischenschritt: $9^3 = 234 \equiv 14 \pmod{55}$, $9^9 = 14^3 \equiv -6 \pmod{55}$, $9^{27} \equiv (-6)^3 = -216 \equiv 4 \pmod{55}$).
			\item \textbf{Signaturprüfung durch Bob (mit $e_A = 3$ und $n_A = 33$):}
			      \[
				      m' = s^{e_A} \bmod n_A = 16^3 \bmod 33 = 4096 \bmod 33 = 4.
			      \]
			      Da $m' = m = 4$, ist die Signatur gültig: Die Nachricht ist geheim übertragen worden, stammt garantiert von Alice und wurde nicht manipuliert.
		\end{enumerate}
	\end{myanswer}
\end{myexercise}